\documentclass[10pt,a4paper]{amsart}
\usepackage[margin=19mm]{geometry}
\usepackage{amsmath,amssymb,mathtools}
\usepackage[scr=rsfs,cal=euler]{mathalfa}
\usepackage{xfrac}
\usepackage[unicode,hidelinks,bookmarks=false]{hyperref}
\newcommand{\ie}{i.e.\ }
\newcommand{\cf}{cf.\ }
\newcommand{\resp}{respectively\ }
\newcommand{\Subsection}{\subsection}
\providecommand{\nobreakdash}{\nobreak\hbox{-}}
\setlength{\emergencystretch}{2em}
\multlinegap0pt
\usepackage{amssymb}
\usepackage{bm}
\usepackage{mathtools}
\usepackage[shortlabels]{enumitem}
\usepackage{xcolor}
\newtheorem{theorem}{Theorem}
\title{Explicit Two-Sided Eigenvalue Bounds for Schr\"{o}dinger Operators with Singular Potentials via Finite Element Method}
\author{Xuefeng Liu}
\date{Source: arXiv:2605.05177v1 (CC BY 4.0)}
\begin{document}
\maketitle

\noindent\emph{Source and licence.} Explicit Two-Sided Eigenvalue Bounds for Schr\"{o}dinger Operators with Singular Potentials via Finite Element Method. Version arXiv:2605.05177v1 (CC BY 4.0), distributed by arXiv under the Creative Commons Attribution 4.0 International licence, http://creativecommons.org/licenses/by/4.0/, as recorded in the arXiv licence metadata for this exact version and shown on the abstract page https://arxiv.org/abs/2605.05177v1. Full licence notice: \texttt{licenses/arxiv-2605.05177-CC-BY-4.0.txt}.\par\smallskip

\noindent\textit{Editorial scope.} Counted unit: the article's theorem thm:graded-eps (source0.tex lines 3225--3240). The article's complete original proof of it is delivered with it (source0.tex lines 3242--3286, one \texttt{proof} environment of 45 lines). No mathematics is written, repaired or re-derived: the statement and the proof are the article's own lines, in the article's own notation, and no retained line is substituted or re-flowed.  Retained uncounted as the source-local context the proof needs: lines 3209--3214; lines 3219--3223.  Nothing was omitted from any retained span, so the package carries no omission record.  No retained line contains a citation, so the wrapper carries no bibliography and no bibliography entry.  No retained line contains a figure environment, an image-inclusion command or drawing code, so no image is delivered and the package declares no asset dependency.  Editorial formatting only, in addition: an AMS \texttt{amsart} preamble that restates the article's own declarations and macros used by the retained lines, namely the environments \texttt{theorem} on separate counters, together with the standard packages those lines need.  The retained environments print in this document as Theorem 1 in order of appearance, whereas the article prints the same items with the numbering of its own full document.  The complete native source of the article as distributed in the arXiv e-print for this exact version is bundled unchanged as \texttt{sources/199930/source0.tex}.
\begin{align}
((V-c_h)u,u)
&=\sum_{K}\int_K (V-V_K)u^2\,dx \notag\\
&=\sum_K \int_K (V-V_K)\bigl(u^2-(u^2)_K\bigr)\,dx,
\label{eq:Vchuident}
\end{align}

\begin{equation}
|((V-c_h)u,u)| \le d_h \|u\|_{H^1(Q)}^2
\qquad \forall\,u\in H^1(Q).
\label{eq:dh-def}
\end{equation}

\begin{theorem}[Direct patch/off-patch estimate, $N=3$]
\label{thm:graded-eps}
Let $N=3$, let $Q\subset \mathbb R^3$ be the auxiliary domain, and let
$V(x)=|x|^{-1}$. Then the constant $d_h$ in \eqref{eq:dh-def} satisfies
\begin{equation}
d_h
\le
\|V-c_h\|_{L^{3/2}(\omega_0)}\,C_6^2
\;+\;
\max_{\{K:\,0\notin \overline K\}}
\frac{h_K}{2}\,\|V-V_K\|_{L^\infty(K)}.
\label{eq:direct-est-c-grd}
\end{equation}
Here $C_6$ is the embedding constant of
$H^1(Q)\hookrightarrow L^6(Q)$.
\end{theorem}

\begin{proof}
By \eqref{eq:Vchuident},
\[
|((V-c_h)u,u)|
\le
\sum_{K\subset \omega_0}\int_K |V-V_K|\,u^2\,dx
\;+\;
\sum_{K\not\subset \omega_0}
\left|\int_K (V-V_K)\bigl(u^2-(u^2)_K\bigr)\,dx\right|.
\]

For the singular patch, Hölder's inequality gives
\[
\sum_{K\subset \omega_0}\int_K |V-V_K|\,u^2\,dx
\le
\|V-c_h\|_{L^{3/2}(\omega_0)} \|u\|_{L^6(\omega_0)}^2
\le
\|V-c_h\|_{L^{3/2}(\omega_0)} C_6^2 \|u\|_{H^1(Q)}^2.
\]

For an off-patch cell $K$, since $K$ is convex,
\[
\|u^2-(u^2)_K\|_{L^1(K)}
\le
\frac{h_K}{2}\|\nabla(u^2)\|_{L^1(K)}
=
h_K\int_K |u||\nabla u|\,dx.
\]
Hence
\[
\|u^2-(u^2)_K\|_{L^1(K)}
\le
h_K \|u\|_{L^2(K)}\|\nabla u\|_{L^2(K)}
\le
\frac{h_K}{2}\bigl(\|u\|_{L^2(K)}^2+\|\nabla u\|_{L^2(K)}^2\bigr).
\]
Therefore
\[
\left|\int_K (V-V_K)\bigl(u^2-(u^2)_K\bigr)\,dx\right|
\le
\frac{h_K}{2}\|V-V_K\|_{L^\infty(K)}
\bigl(\|u\|_{L^2(K)}^2+\|\nabla u\|_{L^2(K)}^2\bigr).
\]
Summing over $K$ with $0\notin \overline K$ yields \eqref{eq:direct-est-c-grd}.
\end{proof}

\end{document}
